\documentclass{csse4400}

% \teachermodetrue

\usepackage{float}
\usepackage{languages}

\title{Terraform Collaboration Guide}
\author{CSSE6400 Teaching Team}
\date{Assignment Guide}

\begin{document}

\maketitle

\aside{
  Use this guide when setting up Terraform for your team assignment.
  Each team should use one shared Terraform state file stored in S3.
}

\section{Overview}

Terraform records the resources it manages in a state file.
When working alone, this is usually a local file named \texttt{terraform.tfstate}.
When working in a team, local state causes problems because each person may have a different view of the infrastructure.

For the assignment, your team should store Terraform state in a shared S3 bucket.
This lets every team member run \texttt{terraform plan} against the same state.
It also allows Terraform to lock the state while someone is applying changes.

\warning{
  Never commit \texttt{terraform.tfstate}, \texttt{terraform.tfstate.backup},
  \texttt{.terraform/}, AWS credentials, or Learner Lab tokens.
}

\section{Choose a Region and Bucket Name}

Your team needs one S3 bucket for Terraform state.
Choose one AWS region and one bucket name before configuring Terraform.
Every team member must use the same region, bucket, and state file path.

\begin{code}[numbers=none]{}
AWS region:       us-east-1
S3 state bucket:  csse6400-team-00-tfstate-a1b2
State file path:  taskoverflow/terraform.tfstate
\end{code}

\noindent
Replace \texttt{us-east-1} with the AWS region your team is using.
Replace \texttt{csse6400-team-00-tfstate-a1b2} with your team's actual bucket name.

\info{
  S3 bucket names are globally unique.
  Include your course, team number, and a short unique suffix in the name.
}

\section{Step 1: Create the Bucket}

One team member should create the bucket.
Do this before running \texttt{terraform init} with the S3 backend.
First set your chosen values in the terminal.

\begin{code}[language=shell,numbers=none]{}
export TF_STATE_BUCKET=csse6400-team-00-tfstate-a1b2
export AWS_REGION=us-east-1
\end{code}

\noindent
For \texttt{us-east-1}, use:

\begin{code}[language=shell,numbers=none]{}
aws s3api create-bucket \
  --bucket "$TF_STATE_BUCKET" \
  --region "$AWS_REGION"

aws s3api put-bucket-versioning \
  --bucket "$TF_STATE_BUCKET" \
  --versioning-configuration Status=Enabled
\end{code}

\noindent
For other regions, change \texttt{AWS\_REGION} and include a location constraint:

\begin{code}[language=shell,numbers=none]{}
export AWS_REGION=ap-southeast-2

aws s3api create-bucket \
  --bucket "$TF_STATE_BUCKET" \
  --region "$AWS_REGION" \
  --create-bucket-configuration LocationConstraint="$AWS_REGION"

aws s3api put-bucket-versioning \
  --bucket "$TF_STATE_BUCKET" \
  --versioning-configuration Status=Enabled
\end{code}

\info{
  Bucket versioning is recommended so that accidental state changes are easier to recover from.
}

\warning{
  The bucket that stores Terraform state should be created before the backend is configured.
  Do not put this same bucket in the Terraform configuration that will use it as a backend.
}

\section{Step 2: Create Backend Files}

In your Terraform directory, create \texttt{backend.tf}.

\begin{code}[language=terraform,numbers=none]{backend.tf}
terraform {
  backend "s3" {}
}
\end{code}

\noindent
Create \texttt{backend.hcl}.
This file tells Terraform which S3 bucket should store the shared state.

\begin{code}[language=terraform,numbers=none]{backend.hcl}
bucket       = "csse6400-team-00-tfstate-a1b2"
key          = "taskoverflow/terraform.tfstate"
region       = "us-east-1"
encrypt      = true
use_lockfile = true
\end{code}

\noindent
The \texttt{region} value must match the region where your team created the bucket.

\info{
  \texttt{use\_lockfile = true} enables S3 state locking.
  This helps stop two team members from applying changes to the same state at the same time.
}

\warning{
  Backend blocks cannot use Terraform variables.
  Do not write \texttt{bucket = var.bucket\_name} inside a backend configuration.
  If Terraform rejects \texttt{use\_lockfile}, update Terraform or ask a tutor.
}

\section{Step 3: Initialise Terraform}

Run Terraform initialisation from the Terraform directory.

\bash{
terraform init -backend-config=backend.hcl
}

\noindent
If this is a fresh project, Terraform should initialise the S3 backend.
If there is already local state, Terraform may ask whether to migrate that state to S3.

\warning{
  Only migrate state from the machine that has the correct current state.
  If multiple people have already run \texttt{terraform apply} locally,
  stop and ask a tutor before migrating.
}

\newpage
\section{Step 4: Ignore Local Files}

Make sure your repository ignores generated Terraform files and local credentials.

\begin{code}[numbers=none]{.gitignore}
.terraform/
terraform.tfstate
terraform.tfstate.backup
*.tfstate
*.tfstate.backup
credentials
aws.env
\end{code}

\noindent
Commit \texttt{backend.tf}.
You may commit \texttt{backend.hcl} if it only contains the team bucket name and no secrets.
If each student needs different local settings,
commit \texttt{backend.example.hcl} instead and keep \texttt{backend.hcl} ignored.

\section{Normal Team Workflow}

Before making changes:

\begin{code}[language=shell,numbers=none]{}
git pull
terraform init -backend-config=backend.hcl
terraform validate
terraform plan
\end{code}

\noindent
Before opening a pull request:

\begin{code}[language=shell,numbers=none]{}
terraform fmt
terraform validate
terraform plan
\end{code}

\noindent
When applying shared infrastructure changes:

\begin{code}[language=shell,numbers=none]{}
git pull
aws sts get-caller-identity
terraform init -backend-config=backend.hcl
terraform plan
terraform apply
\end{code}

\warning{
  Only one person should run \texttt{terraform apply} at a time.
  Tell your team before applying.
  Do not use \texttt{-lock=false} during normal team work.
}

\section{Changing the Backend}

If you change \texttt{backend.tf} or \texttt{backend.hcl},
reinitialise Terraform.
Do this carefully:
a wrong bucket name or wrong region can point your team at a different state file.

\bash{
terraform init -reconfigure -backend-config=backend.hcl
}

\warning{
  If \texttt{terraform plan} wants to destroy something important,
  do not apply until the team understands why.
  Ask a tutor if unsure.
}

\end{document}
